\begin{definition}[The ghost-free system]
  \label{def:aba-ghost-free}
  \lean{PLTS.ABA.Implementation.systemGhostFree}\lean{PLTS.ABA.ABDY.protocol₀}\lean{PLTS.ABA.AFW.protocol₀}\leanok
  \uses{def:aba-implementation, def:aba-protocol, def:aba-afw-protocol}
  The implementation of Definition~\ref{def:aba-implementation} at a one-element ghost record: the same programs, the same network rows and the same coin oracle, with the adversary holding nothing of a round and its two graded-agreement returns free to announce either bit: \leandecl{PLTS.ABA.Implementation.systemGhostFree}.
  The ghost-free systems of the two protocols are that definition at the two implementations, \leandecl{PLTS.ABA.ABDY.protocol₀} and \leandecl{PLTS.ABA.AFW.protocol₀}, each the protocol of Definition~\ref{def:aba-protocol} and of Definition~\ref{def:aba-afw-protocol} with the record dropped and nothing else changed.
  The ghost output of D30 is a relation on the announced bit, which is what puts the three readings in one table: each algorithm instantiates it as the equation between the bit and the record it keeps, and the ghost-free system as the full relation.
\end{definition}
